%=====================================================================
\section{Final --- Diciembre 2025}
%=====================================================================

\begin{ojo}
Este final trae, escrita a mano en la carátula, una nota de puntajes parciales del
alumno rendido (Ej.\,1: 0/3, Ej.\,2: 0,5/3, Ej.\,3: 1/4, nota final 2) pero \textbf{no}
trae desarrollo de respuestas: es prácticamente solo el enunciado. Las resoluciones de
esta sección se elaboraron desde cero contra el material de la materia.
\end{ojo}

% =====================================================================
\subsection{Ejercicio 1 --- Control de entradas por QR con MAC simétrico (Patova SAS)}
% =====================================================================

\begin{consigna}
La empresa Patova SAS ofrece sistemas de control de entradas para eventos (recitales,
partidos, etc.). Uno de sus sistemas utiliza tickets impresos o digitales con un código QR.

El QR de cada entrada contiene la siguiente información: (event\_id, tipo\_de\_entrada,
\#ticket, mac), donde event\_id es un identificador único por evento, tipo\_de\_entrada
indica la zona (palco, vip, campo, etc.), \#ticket es un número único que identifica a
cada ticket, y mac = mac$_{k\_evento}$(event\_id $\|$ \#ticket). Finalmente k\_event es una
clave simétrica de 128 bits definida para cada evento, conocida por el servidor y los
lectores de QR.

En un escenario de mucha concentración de personas es normal que la señal de red funcione
irregularmente, así que los lectores que se utilizan en cada acceso están pensados para
poder operar incluso sin estar conectados al servidor.

El lector decodifica el QR y verifica el mac. Si la verificación es exitosa, revisa que
tipo\_de\_entrada corresponda al lugar de ingreso (por ejemplo, el lector usado en los
ingresos vip solo permite tickets de ese tipo), avisa al servidor por un canal seguro para
confirmar el acceso y si el servidor da el ok permite avanzar. El servidor verifica que
corresponda a un ticket emitido, y que no haya sido utilizado. Si no puede hacerse la
conexión, o si no llega respuesta a tiempo para rechazar usos subsiguientes del mismo
ticket. Una vez retomada la conexión, enviará los números de ticket procesados al servidor.

\begin{enumerate}
  \item[a)] Debido a la autenticación del mac, un atacante no puede inventar números de serie de
  tickets. Explicar qué tipo de fraude si podría realizar abusando de fallas en el control de
  integridad, y explicar cómo evitarlo. (1 pto.)
  \item[b)] Si un atacante consigue un lector y extrae la clave del evento, puede falsificar las
  entradas que quiera. Proponer un cambio para evitar que la extracción de claves del lector no
  permita este tipo de ataque. (1 pto.)
  \item[c)] En un evento hay 4 entradas a diferentes secciones, un lector en cada una. Explicar
  cómo minimizar el problema de que dos personas utilicen la misma entrada en dos puertas
  distintas en un escenario de inestabilidad de la red. (1 pto.)
\end{enumerate}
\end{consigna}

\paragraph{Temas.}
Clase 3 --- MACs y Cifrado Autenticado (\S Qué es un MAC; \S Seguridad de un MAC:
infalsificabilidad Mac-forge; \S método para romper un MAC inseguro cuando no procesa todo
el mensaje; \S CBC-MAC / HMAC); Clase 5 --- Protocolos (\S Nonce y frescura; \S Replay).

\paragraph{Qué necesitás saber.}
\begin{itemize}
  \item \textbf{MAC (Message Authentication Code):} terna $(\mathrm{Gen},\mathrm{Mac},\mathrm{Vrfy})$
  con clave simétrica compartida $k$. Da \textbf{integridad + autenticidad}: si
  $\mathrm{Vrfy}_k(m,t)=1$ el receptor sabe que $m$ fue generado por alguien que conoce $k$
  y no fue alterado en tránsito. \textbf{No impide la modificación} (un atacante en el medio
  puede alterar bits), solo permite \textbf{detectarla}.
  \item \textbf{Infalsificabilidad (Mac-forge):} un MAC es seguro si ningún adversario con
  acceso a un oráculo $f(x)=\mathrm{Mac}_k(x)$ (sin ver $k$) puede producir $(m,t)$ válido con
  $m$ no consultado, salvo con probabilidad despreciable. Como aquí $\mathrm{mac}=
  \mathrm{Mac}_{k\_evento}(\text{event\_id}\|\text{\#ticket})$ con clave de 128 bits, un
  atacante \textbf{no puede} generar un mac válido para un \#ticket inventado sin conocer
  $k\_evento$: no puede fabricar tickets nuevos ni alterar event\_id/\#ticket de uno
  existente sin invalidar el mac.
  \item \textbf{Lo que el MAC NO cubre:} el esquema solo autentica $(\text{event\_id},
  \#\text{ticket})$; no dice nada sobre \emph{cuántas veces} se usa ese par. La detección de
  reutilización (``ticket ya usado'') depende exclusivamente de que el servidor la registre
  y de que los lectores estén sincronizados con él --- y el enunciado aclara que en modo
  degradado (sin conexión) el lector \textbf{no puede rechazar usos subsiguientes del mismo
  ticket} hasta reconectarse.
  \item \textbf{Riesgo de la clave simétrica compartida:} en un MAC (a diferencia de una
  firma digital, Clase 4) \textbf{todos} los verificadores (todos los lectores del evento)
  deben conocer la misma clave $k\_evento$ para poder verificar. Cualquier lector es, por
  diseño, un punto donde la clave está disponible en claro (o debe estarlo para operar);
  extraerla de un lector comprometido da la capacidad de \textbf{generar} macs válidos, no
  solo de verificarlos --- el mismo problema estructural de compartir clave que motiva pasar
  a firma digital (clave privada solo del lado que firma) cuando hay muchos verificadores.
\end{itemize}

\paragraph{Notas y conexiones.}
Este ejercicio aplica directamente la definición de MAC/CBC-MAC/HMAC de la Clase 3 (la
integridad autenticada no impide la modificación, solo la detecta) y el problema de
\emph{replay} de la Clase 5 (reutilizar el mismo mensaje/ticket varias veces mientras no hay
frescura verificada). El punto b) es el mismo argumento por el cual, cuando hay muchos
verificadores independientes, conviene una \textbf{firma digital} (Clase 4): solo quien firma
necesita el secreto ($sk$), todos los demás solo necesitan la clave \textbf{pública} $pk$, así
que comprometer un verificador no permite forjar.

\paragraph{Resolución.}

\textbf{a)} El fraude posible es el \textbf{reuso del mismo ticket (doble uso / ``ticket
compartido'')}: aunque el atacante no puede inventar un \#ticket ni un event\_id nuevos (el
mac se lo impide), \textbf{sí puede reutilizar un QR válido ya emitido} (fotografiarlo,
copiarlo, distribuirlo) porque el mac de $(\text{event\_id}\|\#\text{ticket})$ es siempre el
mismo para ese par: el esquema \textbf{no incluye ningún elemento de frescura} (nonce,
contador de usos, timestamp) que distinga ``primer uso'' de ``uso repetido''. Esto es
particularmente explotable en el \textbf{modo degradado sin conexión}: el enunciado dice
explícitamente que si no llega respuesta del servidor a tiempo \textbf{no se puede rechazar
el uso subsiguiente del mismo ticket}, así que dos (o más) personas con copias del mismo QR
pueden entrar mientras algún lector esté offline. Cómo evitarlo: el control de unicidad no
puede vivir solo en el mac (que es determinístico e idéntico en cada verificación); hay que
que cada lector \textbf{mantenga localmente la lista de \#tickets ya vistos} en ese evento
(incluso offline) y rechace inmediatamente una repetición local, además de sincronizar esa
lista entre lectores apenas se recupera la conexión (tal como el enunciado ya prevé el envío
de los números procesados al reconectar) para cerrar la ventana de doble uso entre lectores
distintos.

\textbf{b)} El problema de fondo es que $k\_evento$ es una \textbf{clave simétrica
compartida por todos los verificadores} (todos los lectores del evento): quien la posee no
solo puede \emph{verificar} macs, sino también \emph{generar} macs válidos para cualquier
$(\text{event\_id}\|\#\text{ticket})$ que quiera, incluidos tickets inexistentes. Extraer la
clave de un único lector comprometido rompe la seguridad de \textbf{todo} el evento.
Propuesta de cambio: reemplazar el esquema por \textbf{firma digital} (Clase 4,
$\mathrm{Sign}_{sk}/\mathrm{Vrfy}_{pk}$): el servidor (o el emisor de tickets) firma
$(\text{event\_id}\|\#\text{ticket})$ con su clave \textbf{privada} $sk$, y cada lector solo
necesita la clave \textbf{pública} $pk$ correspondiente para verificar ($\mathrm{Vrfy}_{pk}$).
Con esto, extraer la clave almacenada en un lector comprometido solo da la clave
\emph{pública} (que ya es pública por definición) y \textbf{no permite falsificar} nuevas
entradas, porque la capacidad de \emph{firmar} queda concentrada del lado del emisor y nunca
viaja a los lectores de campo. Alternativa más acotada sin cambiar de esquema: si se quiere
mantener MAC simétrico, usar HSM en cada lector para que $k\_evento$ nunca salga en claro del
hardware (el lector solo pide la operación de verificación, sin exponer la clave), aunque
esto no elimina el riesgo tan bien como pasar a un esquema asimétrico.

\textbf{c)} Con 4 puertas/secciones y lectores potencialmente desconectados entre sí, el
riesgo es que el mismo \#ticket sea aceptado en dos lectores distintos antes de que se
enteren mutuamente (inestabilidad de red = no hay una única fuente de verdad en tiempo
real). Para minimizarlo:
\begin{itemize}
  \item Igual que en a), cada lector debe llevar su \textbf{propio historial local} de
  \#tickets vistos y, apenas recupere conectividad, \textbf{sincronizar} esa lista con el
  servidor y con los demás lectores (el enunciado ya prevé el envío de los números
  procesados al reconectar), de forma que la ventana de doble uso quede acotada al tiempo
  sin conexión, no indefinida.
  \item Reducir esa ventana: reconectar/sincronizar con la mayor frecuencia posible (o usar
  un canal secundario de baja velocidad pero disponible) en vez de esperar a que vuelva la
  conectividad completa, para minimizar el tiempo en que dos lectores pueden aceptar el mismo
  ticket sin saberlo.
  \item Si tipo\_de\_entrada ya restringe cada \#ticket a una única sección/puerta (como
  ocurre con el ticket vip en el lector vip), un ticket de sección A no debería poder
  aceptarse nunca en el lector de la sección B --- reforzar esa validación de zona en todos
  los lectores reduce a la mitad (o a un cuarto, con 4 puertas) las combinaciones de doble
  uso posibles, aunque no elimina el reuso dentro de la misma puerta.
\end{itemize}
\begin{ojo}
El enunciado no detalla ningún mecanismo de sincronización entre lectores más allá de ``una
vez retomada la conexión enviará los números de ticket procesados al servidor''; la solución
propuesta arriba (historial local + sincronización) es una extensión razonable a partir de
ese texto, pero el material de la materia no describe un protocolo específico de
sincronización distribuida entre lectores de control de acceso, así que esta parte se
resuelve con los principios generales de frescura/anti-replay (Clase 5) aplicados al
escenario, no con una técnica puntual de la materia.
\end{ojo}

% =====================================================================
\subsection{Ejercicio 2 --- SSO con IdP y firma digital}
% =====================================================================

\begin{consigna}
Copérnico SRL posee varias aplicaciones web internas y está diseñando un sistema de login
unificado (SSO), que permita a sus empleados autenticarse una vez y poder utilizar todas las
aplicaciones a las cuales tenga acceso.

Para eso ha decidido instalar un Proveedor de Identidad (IdP) que ofrece una aplicación web
donde los usuarios pueden autenticarse. El IdP posee un par de claves (sk, pk) utilizadas
para realizar firmas digitales. Todas las aplicaciones web conocen la clave pública pk.

Cuando un usuario llega a cualquier aplicación, si no posee una sesión válida, la aplicación
redirige al usuario al IdP, donde se autentica. Cuando el proceso es exitoso (o si el usuario
tiene una sesión ya activa en el IdP), el IdP construye un token T=(id\_usuario, servicio),
donde ``servicio'' es el identificador del servicio que solicitó el login.

Luego calcula S=Sign$_{sk}$(H(T)), donde Sign es una firma digital resistente a
falsificaciones y H(.) es una función de hash libre de colisiones. Finalmente redirige al
usuario a la página del servicio, enviando T y S como parámetros.

Los distintos servicios, ante la presencia de los parámetros T y S, verifican la firma S y si
es válida y el servicio mencionado en el token en el que está validando la firma, crean una
sesión para el usuario.

\begin{enumerate}
  \item[a)] Explicar por qué cada servicio no necesita conocer la contraseña del usuario para
  poder crearle una sesión y eso no tiene un impacto negativo en seguridad. (1 pto.)
  \item[b)] \textquestiondown{}Cómo se puede mejorar el protocolo para evitar el escenario donde un atacante
  consiga de algún modo un par T, S válido e impersone al usuario que los generó en primer
  momento? (1 pto.)
  \item[c)] Un desarrollador propone optimizar el sistema eliminando el campo ``servicio'' del
  token. Desarrollar un ataque que explote esta modificación para ganar acceso a un sistema
  w' robando un par de T,S que permiten acceder al sistema w. Asumir que las comunicaciones
  entre partes están correctamente protegidas por un canal TLS, y explicar un posible
  escenario que le permita ``robar'' esos tokens. (1 pto.)
\end{enumerate}
\end{consigna}

\paragraph{Temas.}
Clase 4 --- Cifrado Asimétrico y Firma Digital (\S Firmas digitales; \S por qué se firma el
hash y no el mensaje; \S infalsificabilidad); Clase 5 --- Protocolos (\S Nonce y frescura;
\S Replay); Segunda Parte, Clase 2 --- Autenticación (\S Autenticación remota SSO/OpenID
Connect).

\paragraph{Qué necesitás saber.}
\begin{itemize}
  \item \textbf{Firma digital:} $\mathrm{Sign}: SK\times P\to S$ con la clave
  \textbf{privada}, $\mathrm{Vrfy}: PK\times P\times S\to\{0,1\}$ con la clave
  \textbf{pública}; da integridad, autenticación y \textbf{no repudio} (solo quien tiene
  $sk$ pudo firmar). Se firma el \textbf{hash} $H(m)$ (libre de colisiones) y no el mensaje
  directamente: entre otras razones, por eficiencia/tamaño fijo de la firma y por evitar
  ataques de falsificación existencial sobre RSA-Signature sin hash.
  \item \textbf{SSO (Single Sign-On):} delegar la autenticación a un sistema externo (el
  IdP) con el cual hay una relación de confianza; conviene cuando ese tercero autentica de
  forma más robusta. Las aplicaciones (``servicios'') confían en el IdP porque pueden
  \textbf{verificar la firma} con la clave pública $pk$, sin necesitar ver jamás la
  contraseña del usuario.
  \item \textbf{Nonce y frescura (Clase 5):} un nonce (número aleatorio de un solo uso)
  permite detectar/evitar \textbf{replays} (reenvío de un mensaje/token viejo). Sin un
  elemento de frescura ni de identidad de destinatario, un par firmado válido puede
  reenviarse o reusarse tal cual en otro contexto.
\end{itemize}

\paragraph{Notas y conexiones.}
El esquema es una aplicación directa de firma digital (Clase 4) al problema de SSO (Segunda
Parte, Clase 2): el IdP firma con su privada, cada servicio verifica con la pública ---
exactamente el patrón ``firmante/verificador'' de la Clase 4, con la ventaja de que el
servicio nunca necesita la contraseña. Las partes b) y c) son, en el fondo, el mismo problema
de \textbf{replay/falta de frescura y de ligazón (binding) del token a un contexto} que
aparece en los protocolos de Needham--Schroeder de la Clase 5: un mensaje firmado/cifrado que
no ata explícitamente todos los datos relevantes (destinatario, tiempo, canal) puede
reutilizarse fuera de su contexto original.

\paragraph{Resolución.}

\textbf{a)} Cada servicio no necesita conocer la contraseña del usuario porque la
\textbf{autenticación queda delegada por completo en el IdP}: el usuario solo ingresa su
contraseña una vez, en el IdP (el único que la conoce/valida), y lo que recibe cada servicio
es un \textbf{token firmado} $T=(\text{id\_usuario},\text{servicio})$ con
$S=\mathrm{Sign}_{sk}(H(T))$. Como la firma es \textbf{resistente a falsificaciones}, ningún
servicio (ni un atacante) puede producir un par $(T,S)$ válido sin conocer $sk$, que solo
tiene el IdP; verificar $S$ con la clave \textbf{pública} $pk$ le alcanza a cada servicio
para confiar en que ese $T$ fue efectivamente emitido por el IdP para ese usuario y ese
servicio. Esto no tiene impacto negativo en seguridad porque: (1) la contraseña nunca se
disemina entre las aplicaciones (menos superficie de ataque: si un servicio se compromete, no
se filtra la contraseña real del usuario, solo se puede leer/verificar tokens ya emitidos);
(2) el modelo de confianza es el mismo que en cualquier esquema firmante/verificador de la
Clase 4: solo el que firma ($sk$, el IdP) puede generar sesiones válidas, y eso es
exactamente lo que se quiere (no repudio/autenticidad respecto del IdP).

\textbf{b)} El riesgo es que un par $(T,S)$ válido, una vez obtenido por un atacante (por
ejemplo interceptándolo o reenviándolo), \textbf{siga siendo válido indefinidamente y en
cualquier contexto}: no hay nada en el esquema que ate ese par a un momento (frescura) ni a
un canal/sesión concretos, de modo que reenviarlo (\emph{replay}) alcanza para impersonar.
Mejora: incorporar al token $T$ (antes de firmarlo) un elemento de \textbf{frescura tipo
nonce/timestamp} (Clase 5), por ejemplo $T=(\text{id\_usuario},\text{servicio},\text{nonce o
timestamp de emisión})$, y que cada servicio (a) verifique que el token no está
\textbf{expirado} (ventana de validez corta) y (b) \textbf{registre/invalide} el nonce ya
usado, de forma que un mismo $(T,S)$ no pueda reutilizarse dos veces ni fuera de su ventana
temporal. Esto convierte el token en un ``desafío-respuesta con vencimiento'' en vez de un
secreto reusable indefinidamente, igual que los nonces $r_1,r_2$ evitan el replay de la clave
de sesión en Needham--Schroeder.

\textbf{c)} Si se elimina el campo \texttt{servicio} del token, éste queda como
$T=(\text{id\_usuario})$, y $S=\mathrm{Sign}_{sk}(H(T))$ ya \textbf{no ata el token a ningún
servicio en particular}: la firma certifica ``el IdP autenticó a id\_usuario'', pero no
certifica \emph{para qué servicio}. Ataque: un atacante que controla (o compromete) el
sistema $w'$ hace que un usuario legítimo inicie sesión en $w'$ vía el mismo IdP; el IdP
genera un $(T,S)$ válido (firma legítima del IdP para ese id\_usuario) que $w'$ recibe como
parte del flujo normal de redirección. Como el canal TLS protege la confidencialidad e
integridad \emph{en tránsito}, pero \textbf{no impide que el propio destino $w'$, que es
quien recibe $T,S$ en la práctica, los reutilice él mismo}: el operador (o código malicioso)
de $w'$ toma ese mismo $(T,S)$ --- que ya no menciona ningún servicio destino --- y lo
reenvía como parámetros a $w$ (el sistema real que el atacante quiere atacar). Como $w$ solo
verifica $\mathrm{Vrfy}_{pk}(H(T),S)=1$ y ya no puede comparar ningún campo ``servicio''
contra sí mismo (porque ese campo no existe), acepta el token y crea una sesión para
id\_usuario en $w$, aun cuando el usuario nunca intentó autenticarse contra $w$. En síntesis:
sin el campo \texttt{servicio} el token deja de estar \textbf{ligado a su destino} y
cualquier servicio que reciba honestamente un $(T,S)$ puede reutilizarlo (``confused deputy''
entre servicios) para autenticarse ante \emph{cualquier otro} servicio que confíe en la misma
$pk$ --- el TLS protege el transporte, no la reutilización del mismo mensaje firmado en un
contexto distinto para el cual no fue emitido.

% =====================================================================
\subsection{Ejercicio 3 --- Hipótesis de falla para la app CriptoTotal}
% =====================================================================

\begin{consigna}
La empresa CriptoTotal ofrece una aplicación web y móvil para que los usuarios puedan ver,
en un solo lugar, sus inversiones en distintas criptomonedas realizadas en múltiples
exchanges.

El registro se hace con email y contraseña, con opción de activar un segundo factor de
autenticación que envía códigos de uso único por SMS.

El usuario puede vincular cuentas de distintos exchanges (por ejemplo Trinance, Saken,
Rupio) cargando sus API keys y secret keys. La documentación comercial afirma que ``las
claves se almacenan cifradas en nuestra base de datos''.

Periódicamente, un servicio backend consulta las APIs de los exchanges vinculados,
descarga balances y movimientos y los guarda en una base de datos relacional.

La aplicación utiliza tokens JWT para manejar la sesión del usuario. El frontend (Single
page app en JavaScript, que corre en el navegador de cada usuario) almacena el JWT en
localStorage y lo envía en el header Authorization: Bearer\ldots en cada request API REST.

Existe un panel de administración al que acceden empleados de soporte para ver el estado de
las cuentas, saldos agregados, últimos accesos y errores de sincronización.

La versión ``premium'' de la app se paga por suscripción mensual a través de un proveedor de
pagos -- la app redirige al proveedor y este último controla todo el flujo de checkout,
captura de información de pago y ejecución. Al confirmarse el pago, se marca en una tabla la
cuenta del usuario como premium = true y se habilitan funciones adicionales (reportes
históricos, alertas por precio, etc.).

En función de la descripción del sistema, establecer 4 hipótesis de falla plausibles
siguiendo la metodología de pentesting. Por cada hipótesis definir una prueba siguiendo los
lineamientos de la metodología. (1 pto. c/u.)

Aclaración: No utilizar problemas genéricos como hipótesis de falla. Plantear situaciones y
lugares concretos donde podría estar el problema. Por ejemplo, en lugar de decir ``puede
haber un buffer overflow'', decir ``es probable que haya un buffer overflow en la lectura del
parámetro x que debe interpretarse como número y puede tener un buffer pequeño dado que los
números esperados tienen menos de 4 dígitos''.
\end{consigna}

\begin{ojo}
El PDF del examen trae, escritas a mano al lado de cada párrafo del enunciado, anotaciones
breves (``Sensitive Data Exposure'', ``Exposure/Improper authentication'', ``DoS'',
``Missing auth'', ``Interception''). No forman parte del enunciado tipeado del examen ni
están desarrolladas: parecen una guía de corrección del docente, no una respuesta del
alumno. Se las tuvo en cuenta como pista de \emph{dónde} mirar, pero las 4 hipótesis de abajo
se construyen y justifican exclusivamente con los conceptos y terminología del material de la
materia (no se usa terminología externa como ``OWASP API Top 10'' salvo que ya figure en los
resúmenes).
\end{ojo}

\paragraph{Temas.}
Clase 6 --- Pentesting (\S Metodología de Hipótesis de Falla completa: Recolección $\to$
Hipótesis $\to$ Prueba $\to$ Generalización $\to$ Eliminación); Segunda Parte, Clase 2 ---
Autenticación (\S Multifactor y multicanal: independencia de factores y entrelazamiento);
Clase 3 --- Principios de Diseño y Vulnerabilidades (\S XSS CWE-79; \S Mínimo privilegio /
separación de privilegios); Clase 7 --- Protección de Datos (\S Encripción at-rest: HSM,
Keystore, esquema KEK/DEK).

\paragraph{Qué necesitás saber.}
\begin{itemize}
  \item \textbf{Metodología de hipótesis de falla (memorizar el orden):} Recolección de
  información $\to$ \textbf{Planteo de hipótesis de falla} $\to$ Prueba $\to$ Generalización
  $\to$ Eliminación (opcional). El paso de \textbf{Hipótesis} exige suponer la existencia de
  una vulnerabilidad concreta en una parte concreta del sistema (no genérica). El paso de
  \textbf{Prueba} debe diseñarse lo menos intrusivo posible, ser \textbf{repetible y
  concluyente}, y hoy se asume que hay que \textbf{intentar explotar} la vulnerabilidad para
  demostrarla (salvo que analizar documentación/comportamiento ya alcance).
  \item \textbf{Independencia de factores / entrelazamiento (2FA):} el multifactor es más
  seguro porque obliga a comprometer \textbf{varios canales independientes}; si los factores
  quedan \textbf{entrelazados} (p.\ ej.\ poder recuperar/reiniciar el primer factor usando el
  segundo canal) se pierde esa independencia --- ``dos llaves para la misma cerradura''.
  \item \textbf{XSS (CWE-79):} inyección de código JS que termina \textbf{ejecutado por el
  navegador de la víctima}, aprovechando la confianza del sitio; permite \textbf{robar la
  sesión}. Un JWT guardado en \texttt{localStorage} es accesible por \emph{cualquier} script
  que corra en esa página, a diferencia de una cookie con flag \texttt{HttpOnly} (no
  mencionado explícitamente en el material, pero la propiedad de que localStorage es legible
  por JS sí lo está vía la definición de XSS).
  \item \textbf{Mínimo privilegio / separación de privilegios:} dar a cada sujeto solo los
  permisos mínimos necesarios (``por default, no tener acceso a nada''); una tarea crítica no
  debería quedar bajo un solo sujeto sin controles.
  \item \textbf{Encripción at-rest y esquema KEK/DEK:} para proteger secretos guardados
  (como las API keys/secret keys de los exchanges) se recomienda \textbf{HSM} (opera sin
  exponer la clave) o, en software, un \textbf{Keystore} con clave máster, o directamente el
  esquema por capas \textbf{KEK/DEK} (una \emph{Key Encryption Key} cifra la \emph{Data
  Encryption Key}, que a su vez cifra los datos). Decir simplemente ``las claves se
  almacenan cifradas'' no especifica si existe esa separación de capas ni dónde vive la
  clave que las descifra.
\end{itemize}

\paragraph{Notas y conexiones.}
Es la aplicación directa de la metodología de hipótesis de falla (Clase 6) a un sistema
descrito en detalle, con la restricción explícita del enunciado de no usar hipótesis
genéricas. Cada hipótesis abajo señala un lugar concreto (una línea del enunciado) y una
prueba puntual, repetible y poco intrusiva, siguiendo el Paso 3 de la metodología.

\paragraph{Resolución.}

\textbf{Hipótesis 1 --- Entrelazamiento del segundo factor SMS con la recuperación de
contraseña.} Lugar concreto: el flujo de \emph{recuperación de cuenta} de CriptoTotal, que
probablemente usa el mismo canal (email o el propio SMS) para resetear la contraseña. Es
probable que, si el usuario activó 2FA por SMS, el flujo de ``olvidé mi contraseña'' permita
recuperar el acceso con \textbf{solo el código SMS} (sin la contraseña original), lo cual
\textbf{entrelaza} el primer factor (contraseña) con el segundo (SMS): quien controla el
número de teléfono (p.\ ej.\ mediante SIM swapping o intercepción del SMS) termina teniendo
``dos llaves para la misma cerradura'' y puede tomar la cuenta completa sin conocer nunca la
contraseña real. Prueba: registrar una cuenta de prueba con 2FA por SMS activo, iniciar el
flujo de recuperación de contraseña y verificar si el sistema permite fijar una contraseña
nueva usando únicamente el código recibido por SMS (sin pedir ningún dato adicional
independiente del canal SMS); si es así, queda demostrado el entrelazamiento entre ambos
factores.

\textbf{Hipótesis 2 --- Ausencia de capas KEK/DEK en el cifrado de las API keys/secret keys
de los exchanges.} Lugar concreto: la frase de la documentación comercial ``las claves se
almacenan cifradas en nuestra base de datos'', sin especificar dónde vive la clave que
descifra esas API keys. Es probable que exista una \textbf{única clave máster} (o ninguna
separación KEK/DEK) que cifra directamente todas las secret keys de todos los usuarios, y que
esa clave máster resida accesible desde el mismo backend que hace las consultas periódicas a
los exchanges (mismo proceso/config/variable de entorno); si esa clave máster se compromete
(o se filtra junto con la base de datos, por ejemplo en un backup), \textbf{todas} las secret
keys de todos los usuarios de todos los exchanges quedan expuestas de una sola vez. Prueba:
en un ambiente de prueba, revisar la configuración/infraestructura del servicio backend que
descifra las secret keys para consultar los exchanges y verificar si existe una separación
KEK/DEK (rotación de la clave que cifra sin re-cifrar los datos) o un HSM/Keystore, o si en
cambio hay una única clave estática embebida junto al mismo servicio que accede a la base de
datos --- si la misma credencial que abre la base de datos alcanza para descifrar las secret
keys, la hipótesis queda confirmada.

\textbf{Hipótesis 3 --- Robo del JWT de sesión vía XSS por estar en localStorage.} Lugar
concreto: el frontend (SPA en JavaScript) guarda el JWT en \texttt{localStorage} y lo reenvía
en el header \texttt{Authorization: Bearer\ldots}. Es probable que exista algún punto de
entrada de datos no sanitizado en la SPA (por ejemplo, un campo donde se muestra el nombre de
un exchange vinculado, un comentario o cualquier dato reflejado en la interfaz sin escapar)
que permita inyectar JavaScript (XSS, CWE-79); dado que \texttt{localStorage} es legible por
\emph{cualquier} script que corra en el origen de la página, un XSS exitoso permite leer el
JWT y exfiltrarlo, dando control completo de la sesión (incluyendo ver inversiones y
potencialmente operar sobre cuentas vinculadas) sin necesitar la contraseña ni el 2FA.
Prueba: identificar un campo de entrada de usuario que se refleje en la UI de la SPA (nombre
de cuenta, alias de un exchange vinculado, etc.), intentar inyectar una carga XSS de prueba
inocua (por ejemplo que solo lea \texttt{localStorage.getItem('token')} y lo muestre en un
\texttt{alert} o lo envíe a un endpoint de control propio) y verificar si se ejecuta en el
navegador; si el script corre y logra leer el JWT, la hipótesis queda confirmada de forma
repetible.

\textbf{Hipótesis 4 --- Panel de administración de soporte sin mínimo privilegio (acceso a
saldos agregados de cualquier usuario).} Lugar concreto: ``existe un panel de administración
al que acceden empleados de soporte para ver el estado de las cuentas, saldos agregados,
últimos accesos y errores de sincronización''. Es probable que ese panel no aplique
\textbf{mínimo privilegio}: en lugar de que cada agente de soporte solo pueda ver los datos
del/de los ticket(s) de soporte que tiene asignados, es probable que cualquier cuenta de
soporte pueda consultar saldos agregados y datos de \emph{cualquier} usuario del sistema sin
una necesidad de negocio puntual ni un registro/auditoría de qué usuario consultó a qué otro
usuario, violando el principio de dar a cada sujeto ``solo los permisos mínimos necesarios''.
Prueba: con una cuenta de soporte de prueba con permisos mínimos asignados, intentar consultar
en el panel los saldos agregados y últimos accesos de un usuario \textbf{no relacionado} con
ningún caso de soporte abierto por esa cuenta (por ejemplo cambiando el identificador de
usuario en la URL o parámetro de la consulta del panel) y verificar si el panel entrega esos
datos sin restricción ni quedar registrado como acceso fuera de alcance; si el acceso se
concede igual, la hipótesis queda confirmada.

\begin{ojo}
El enunciado también describe el flujo de pago premium delegado a un proveedor externo
(``el proveedor de pagos controla todo el flujo de checkout... al confirmarse el pago se
marca premium=true''), que sugiere una posible quinta hipótesis (validar cómo se confirma
el pago hacia la propia app y si ese aviso podría falsificarse). No se desarrolla como una
de las 4 hipótesis pedidas porque el material de la materia no cubre en detalle los mecanismos
de confirmación de pagos de un proveedor externo (webhooks firmados, callbacks, etc.); se
señala para que el usuario decida si quiere agregarla igualmente o dejarla fuera por no estar
sustentada en el material.
\end{ojo}
